%----------------------------------------------------------------------------------------
% NOTATIE EN FIGUREN VAN DE CURSUS PREDICATENLOGICA
% Staat in de preambule en niet in een hoofdstuk: ook de lijsten met oplossingen gebruiken dit.
%----------------------------------------------------------------------------------------
% Bewaarde grafiek, geplaatst door maak-pdf.py: bestand en pagina in figuren-cache/
\newcommand{\grafiekbestand}[2]{\includegraphics[page=#2,trim=3bp 3bp 3bp 3bp]{figuren-cache/#1}}
\usetikzlibrary{positioning}

% Zoals stappen, maar met een bredere kolom voor de berekening: voor lange formules
\newenvironment{stappenbreed}{%
  \par\renewcommand{\arraystretch}{1.25}%
  \noindent\begin{tabular}{@{}>{$\displaystyle}p{9.8cm}<{$}@{\hspace{4mm}}>{\RaggedRight\small\itshape\arraybackslash}p{4.9cm}@{}}%
}{\end{tabular}\par}

% Zoals stappen, met eigen kolombreedtes: \begin{stappenvrij}{6.2cm}{3.8cm}. Voor naast een figuur.
\newenvironment{stappenvrij}[2]{%
  \par\renewcommand{\arraystretch}{1.25}%
  \noindent\begin{tabular}{@{}>{$\displaystyle}p{#1}<{$}@{\hspace{3mm}}>{\RaggedRight\small\itshape\arraybackslash}p{#2}@{}}%
}{\end{tabular}\par}

% Tekst of berekening links, figuur rechts: \naast[breedte rechts]{links}{rechts}
\newsavebox{\naastbox}
\newcommand{\naast}[3][]{%
  \par\noindent\sbox{\naastbox}{#3}%
  \begin{minipage}[t]{\dimexpr\linewidth-\wd\naastbox-3mm\relax}\vspace{0pt}#2\end{minipage}\hfill
  \begin{minipage}[t]{\wd\naastbox}\vspace{0pt}\usebox{\naastbox}\end{minipage}\par\addvspace{1mm}}
% Een rij figuren, gecentreerd
\newcommand{\figuren}[1]{\par\addvspace{1mm}{\centering #1\par}\addvspace{1.5mm}}
% Opgave met een grafiek rechts: \begin{metgrafiek}[breedte]{tikz} opgave en deelvragen \end{metgrafiek}
\newsavebox{\grafiekbox}
\newenvironment{metgrafiek}[2][]{%
  \sbox{\grafiekbox}{#2}%
  \begin{minipage}[t]{\dimexpr\linewidth-\wd\grafiekbox-3mm\relax}%
}{%
  \end{minipage}\hfill
  \begin{minipage}[t]{\wd\grafiekbox}\vspace{-2mm}\usebox{\grafiekbox}\end{minipage}\par}

% Rooster met assen: \rooster{xmin}{xstap}{xmax}{ymin}{ystap}{ymax}, met getallen op de assen (gehele waarden).
% \roosterleeg doet hetzelfde zonder getallen.
\newcommand{\roosterleeg}[6]{%
  \draw[black!14,very thin,xstep=#2,ystep=#5] (#1,#4) grid (#3,#6);
  \draw[->,grijs] (#1,0) -- ({#3+.4*#2},0) node[right,text=black] {$x$};
  \draw[->,grijs] (0,#4) -- (0,{#6+.4*#5}) node[above,text=black] {$y$};}
\newcommand{\rooster}[6]{%
  \roosterleeg{#1}{#2}{#3}{#4}{#5}{#6}
  \pgfmathtruncatemacro{\rtweex}{#1+#2}\pgfmathtruncatemacro{\rtweey}{#4+#5}
  \foreach \t in {#1,\rtweex,...,#3}{\ifnum\t=0\else\node[below,inner sep=1.5pt,text=black] at (\t,0) {$\t$};\fi}
  \foreach \t in {#4,\rtweey,...,#6}{\ifnum\t=0\else\node[left,inner sep=1.5pt,text=black] at (0,\t) {$\t$};\fi}}

% Getallenverzamelingen
\newcommand{\N}{\mathbb{N}}
\newcommand{\Z}{\mathbb{Z}}
\newcommand{\Q}{\mathbb{Q}}
\newcommand{\R}{\mathbb{R}}
\newcommand{\C}{\mathbb{C}}
% Logische symbolen, zoals in Logica en bewijzen: \en, \of, \niet, \dan (implicatie), \asa (equivalentie)
\newcommand{\en}{\wedge}
\newcommand{\of}{\vee}
\newcommand{\niet}{\neg}
\newcommand{\dan}{\Rightarrow}
\newcommand{\asa}{\Leftrightarrow}
% Deelbaarheid: a \deelt b
\newcommand{\deelt}{\mid}
% Een predicaat of constante in rechte letters als het een woord is: \pr{Kent}
\newcommand{\pr}[1]{\mathrm{#1}}

% FIGUREN, binnen een tikzpicture
% Venndiagram in de kleur van het hoofdstuk
\tikzset{verz/.style={draw=accent,thick,fill=accent,fill opacity=.10,text opacity=1}}
% Een relatie op een eindig universum als pijlen: \node[elt] (a) at (0,0) {$a$}; \draw[pijl] (a) -- (b);
% Een pijl van x naar y betekent: R(x, y) is waar. Een lus: \draw[pijl] (a) to[lus] (a);
\tikzset{
  elt/.style={circle,draw=accent,thick,fill=accent!10,minimum size=6.2mm,inner sep=0pt,font=\small},
  eltfel/.style={elt,draw=fel!85!black,fill=fel!30},
  pijl/.style={-{Stealth[length=2.1mm]},thick,accent,shorten >=.6pt,shorten <=.6pt},
  pijlfel/.style={pijl,fel!85!black},
  lus/.style={out=120,in=60,looseness=6,min distance=5mm},
  lusonder/.style={out=-60,in=-120,looseness=6,min distance=5mm},
  lusrechts/.style={out=30,in=-30,looseness=6,min distance=5mm},
  luslinks/.style={out=150,in=210,looseness=6,min distance=5mm},
  gebiedstijl/.style={fill=accent!22},
}
